\begin{frame}{Le français louisianais}
  \only<1>{
    \footnotesize
    \begin{itemize}
      \item Still French
      \begin{itemize}
        \item[$\to$] Not inferior, simpler, nor a different language
      \end{itemize}
      \item Vocabulary
      \begin{itemize}
        \item[$\to$] Mostly shared with French, sometimes different meanings, sometimes unique words, sometimes English words
      \end{itemize}
      \begin{itemize}
        \item[] (don't say \lexi{catin} in France)
      \end{itemize}
      \item Grammar
      \begin{itemize}
        \item[$\to$] Mostly shared with French, some things more complicated, some things simpler
      \end{itemize}
      \item Pronunciation
      \begin{itemize}
        \item[$\to$] Like the difference between the English of someone from England and the English of someone from the rural Deep South
      \end{itemize}
    \end{itemize}
  }
  \only<2->{
    \begin{center}
      There are many ways to speak Louisiana French!
    \end{center}
  }
\end{frame}